% Grothendieck's « problem three » (Buffalo, 2 July 1973): generation under limits and algebraic theories relative to a class of limits — standalone article for Hexagon.
% Grew out of Appendix B of docs/buffalo-open/article.tex. Quotations from the recordings are kept to a few words.
% pdfLaTeX. Compile: tectonic -X compile article.tex.
\documentclass[11pt,a4paper]{amsart}
\usepackage[T1]{fontenc}
\usepackage[utf8]{inputenc}
\usepackage{lmodern,microtype,amssymb,mathtools}
\usepackage[margin=2.7cm]{geometry}
\usepackage[hidelinks]{hyperref}

\newtheorem{theorem}{Theorem}[section]
\newtheorem{proposition}[theorem]{Proposition}
\newtheorem{lemma}[theorem]{Lemma}
\newtheorem{corollary}[theorem]{Corollary}
\theoremstyle{definition}
\newtheorem{definition}[theorem]{Definition}
\newtheorem{question}[theorem]{Question}
\theoremstyle{remark}
\newtheorem{remark}[theorem]{Remark}

\newcommand{\q}[1]{\guillemotleft\,#1\,\guillemotright}
\newcommand{\ts}[3]{\href{https://www.youtube.com/watch?v=#1&t=#2}{#3}}
\newcommand{\Set}{\mathbf{Set}}
\newcommand{\Cat}{\mathbf{Cat}}
\newcommand{\Hom}{\operatorname{Hom}}

\title[Grothendieck's problem three]{Grothendieck's \q{problem three}:\\ generation under limits and algebraic theories\\ relative to a class of limits}
\author{Michel Hua}
\address{Commutator, Paris, France}
\email{michel@commutator.io}
\thanks{ORCID 0009-0002-8304-0683.}
\date{8 November 2026}
\subjclass[2020]{18C10; 18A35, 18C30, 18C35, 18F10, 01A60}
\keywords{algebraic theories, closure under limits, locally presentable categories, sketches, Grothendieck topoi, Buffalo 1973, Grothendieck}

\begin{document}
\begin{abstract}
In a lecture on topos theory given at Buffalo on 2 July 1973, Grothendieck asked what he called \q{problem three}. Fix a class $D$ of limits and colimits, and call $I\to R$ basic when $R$ is reached from a set $I$ of objects by a small transfinite sequence of free constructions: the free $D$-envelope, localisations, and the adjunction of arrows. Is $I\to R$ basic exactly when $R$ is the closure of the image of $I$ under $D$-limits and $D$-colimits? He expected one implication and hoped for the other. Assuming, as he did, that the free constructions exist, we prove that basic implies generation for every class $D$, and that $I\to R$ is basic if and only if $R$ is generated by $I$ and is algebraic over some other base. For a small class $D$ the two notions coincide, with the base fixed. For Grothendieck topoi they coincide as well, through a lemma on closures in a topos; two related statements of his course on topoi are false as stated. For all small limits the answer is negative: every set is an equaliser of two maps between powers of $2$, so that $\{2\}$ generates $\Set$ in two steps, but $\Set$ is not algebraic, since $2$ is not presentable in $\Set^{\mathrm{op}}$. The parts of this counterexample about sets are checked in Lean. An example of Adámek and Rosický, which needs a proper class of steps, gives a second counterexample and shows that the number of steps is not bounded in general, although it is bounded for algebraic categories. The converse for mixed large classes, such as finite limits together with all colimits, stays open.
\end{abstract}
\maketitle

\section{The question}\label{sec:question}

From April to July 1973 Grothendieck taught a course on topos theory at the State University of New York at Buffalo. Its recordings are online, on a YouTube channel under the pseudonym \q{Erwin Schulhoff}, and the edition of which this note is part has indexed them and reconstructed the blackboards \cite{repo}. The course spends several sessions on algebraic theories relative to a class of limits. One fixes a class $D$ of diagram shapes, considers the 2-category $\Delta=\Cat_D$ of categories with $D$-limits and $D$-colimits, and builds theories from a set $I$ of base objects by free constructions in $\Delta$: the free $D$-envelope of $I$, localisations, equalisations, and the adjunction of arrows. A theory obtained in this way is \emph{basic} over $I$, and an object of $\Delta$ that is basic over some $I$ is \emph{algebraic}; Section~\ref{sec:framework} gives the precise definitions.

On 2 July 1973 (video \texttt{4gynb9uTEE4}, from \ts{4gynb9uTEE4}{2h3m54s}{2:03:54}) he asked whether there is a simpler way to say that $I\to R$ is basic. His first hope, that it amounts to faithfulness of $\Hom_\Delta(R,C)\to C^I$ for every $C$, he set aside at once, citing a counterexample that he attributes to Isbell (the name is uncertain in the recording). He then proposed a second characterisation, which he numbered problem three. Assume that the free constructions exist (his \q{problems one and two}). Is $I\to R$ basic if and only if every strictly full subcategory of $R$ that contains the image of $I$ and is closed under $D$-limits and $D$-colimits is all of $R$? He reformulated this condition as a transfinite saturation of the image of $I$ under $D$-limits and $D$-colimits, and called the question \q{the main problem which remains open} in this generality; he had proved neither direction, expected the forward one, and hoped for the converse. He then remarked (\ts{4gynb9uTEE4}{2h8m14s}{2:08:14}) that for a small class $D$ the question should reduce to small categories, and conjectured (\ts{4gynb9uTEE4}{2h10m23s}{2:10:23}) that for an algebraic $R$ the base $I$ can be enlarged so that only axioms, and no new arrows, remain to be imposed.

On 10 July 1973 (same video, from \ts{4gynb9uTEE4}{3h37m36s}{3:37:36}) he retracted several claims made on 2 July: that strict cogeneration is transitive, that the image of $I$ stays strictly cogenerating at each step of the construction, that a characterisation by a small strictly cogenerating subcategory holds, and that no transfinite construction is needed. As a counterexample to the last characterisation he gave $R=\Set$ (\ts{4gynb9uTEE4}{3h43m0s}{3:43:00}, \ts{4gynb9uTEE4}{3h44m8s}{3:44:08}). The present note answers problem three as he finally posed it. It grew out of an appendix of a survey of the open problems of the Buffalo lectures \cite{open}, which it replaces. The transcripts we used are first passes from captions that no one has checked.

\subsection*{Results}
Write (G) for the condition of problem three. Assuming the free constructions exist (hypothesis (H) of Section~\ref{sec:framework}), we prove:
\begin{enumerate}
\item basic implies (G), for every class $D$, small or large (Theorem~\ref{t:basicG});
\item $I\to R$ is basic if and only if (G) holds and $R$ is algebraic over some base (Theorem~\ref{t:reduction}); so problem three, for a given $\Delta$, asks exactly whether (G) for some small $I$ implies that $R$ is algebraic (Corollary~\ref{c:reduction});
\item for a small class $D$, basic and (G) are equivalent with the base fixed (Theorem~\ref{t:small});
\item for Grothendieck topoi they are equivalent too (Theorem~\ref{t:topos}), through the fact that the closure of a small set of objects of a topos under finite limits and small colimits is again a topos (Lemma~\ref{l:topos}); two related statements of the course on topoi are false as stated (Proposition~\ref{p:toposfalse});
\item for $D$ all small limits the converse fails: $\{2\}\to\Set$ satisfies (G) after exactly two steps of saturation (Proposition~\ref{p:eq} and Corollary~\ref{c:twosteps}), but $\Set$ is not algebraic (Theorem~\ref{t:notalg}); the facts about sets are checked in Lean (Remark~\ref{r:lean});
\item in an algebraic category the saturation of any generating set stops after a small number of steps (Proposition~\ref{p:bound}), while a cocomplete category given by Adámek and Rosický is generated by one object in a proper class of steps; so it is a second counterexample, and the number of steps is unbounded in general (Corollary~\ref{c:ar}).
\end{enumerate}
What stays open is the converse for mixed large classes, such as finite limits with all small colimits but no exactness condition (Section~\ref{sec:open}).

\subsection*{What is in print}
Our search, recorded in the edition's repository \cite{repo}, found the following, which we use and credit as we go. The transfinite description of the closure of a full subcategory under a class of limits, and the bound on the number of steps when the shapes are small, are in Kelly's book \cite[\S3.5]{kellybook}. Kelly also proves that the free cocompletion of a category with respect to a class of colimits is the closure of the representables under those colimits \cite[Thm~5.35]{kellybook}, which is the first step of Theorem~\ref{t:basicG}, and he gives the sketch theory that we use to show that $\Set^{\mathrm{op}}$ is not algebraic \cite[ch.~6]{kellybook}. That every set is an equaliser of two maps between powers of $2$ is a form of the duality between sets and complete atomic Boolean algebras; it also follows from Paré's theorem that the power-object functor of a topos is monadic \cite{pare}. That $\Set^{\mathrm{op}}$ is not locally presentable goes back to Gabriel and Ulmer \cite{gu}; that algebraic categories, in the present sense and up to duality, are the locally presentable ones is the content of Gabriel--Ulmer duality \cite{gu}, \cite[ch.~1]{ar}. The example of a closure under colimits that needs a proper class of steps is due to Adámek and Rosický \cite[Example~6.38]{ar}; we learnt of it from a MathOverflow discussion \cite{mo}. Faithfulness and density in Isbell's sense \cite{isbell} are related but different notions (Remark~\ref{r:faith}). We found no statement in print of Theorems~\ref{t:basicG}, \ref{t:reduction}, \ref{t:small} or~\ref{t:topos}, nor of Lemma~\ref{l:topos}, nor of the use of these examples against problem three. Most of the arguments are routine for specialists; we give them in full so that the note can be read on its own.

\section{The framework}\label{sec:framework}

We work in ZFC with one Grothendieck universe; \q{small} means in the universe, and categories are locally small unless they are said to be large.

\subsection{The 2-category $\Delta$}
Fix $D=(D',D'')$, two sets of small diagram shapes. The 2-category $\Delta=\Cat_D$ has as objects the locally small categories with all limits of shape in $D'$ and all colimits of shape in $D''$; as 1-cells the \emph{$D$-functors}, which preserve these limits and colimits; and as 2-cells all natural transformations. Write $\Hom_D(R,C)$ for the category of $D$-functors $R\to C$ and natural transformations. The class $D$ is \emph{small} if $D'$ and $D''$ are small sets of small shapes. Examples are finite limits ($D'$ the finite categories, $D''=\varnothing$), all small limits ($D'$ all small categories), all small colimits, and the class of topoi (finite limits and all small colimits, together with exactness conditions; Section~\ref{sec:topos}).

\subsection{Hypothesis (H)}
Grothendieck assumed that certain free constructions exist in $\Delta$; he called this \q{problems one and two solved} (\ts{4gynb9uTEE4}{2h4m30s}{2:04:30}). Each construction is defined by a 2-universal property against every $C\in\Delta$.

\begin{definition}[Hypothesis (H)]\label{d:H}
The following constructions exist in $\Delta$.
\begin{enumerate}
\item The \emph{envelope} $j\colon I\to\widetilde I$ of a small category $I$: restriction along $j$ is an equivalence $\Hom_D(\widetilde I,C)\to\Hom(I,C)$.
\item The \emph{localisation} $R\to R[\Sigma^{-1}]$ at a small set $\Sigma$ of arrows of $R$: restriction from $\Hom_D(R[\Sigma^{-1}],C)$ to $\Hom_D(R,C)$ is fully faithful, with essential image the $D$-functors that invert $\Sigma$.
\item The \emph{equalisation} of a small set of parallel pairs of arrows of $R$, defined in the same way.
\item The \emph{adjunction of arrows}: for a small set $H$ of pairs $(a_h,b_h)$ of objects of $R$ (his rule~3, \ts{PqCQM3Pcl6Q}{8h20m37s}{8:20:37}), a $D$-functor $R\to R\langle H\rangle$ with
\[\Hom_D(R\langle H\rangle,C)\simeq\{(\psi,\beta):\psi\in\Hom_D(R,C),\ \beta_h\colon\psi a_h\to\psi b_h\text{ for }h\in H\},\]
where a morphism $(\psi,\beta)\to(\psi',\beta')$ is a transformation $\psi\to\psi'$ compatible with the $\beta$'s.
\item 2-colimits of small chains, and 2-pushouts (pushouts are used only in Section~\ref{sec:reduction}).
\end{enumerate}
\end{definition}

For $D'=\varnothing$ or $D''=\varnothing$ and $D$ small, Kelly's book provides envelopes and localisations as small categories \cite[Thm~5.35, Prop.~6.21, Cor.~6.22, Thm~6.23]{kellybook}; for topoi the constructions are classifying topoi, 2-pushouts and subtopoi. We have not found a reference for mixed small $D$, with $D'$ and $D''$ both nonempty; Grothendieck asserts it without proof (\ts{PqCQM3Pcl6Q}{5h43m33s}{5:43:33}). Statements below that use (H) say so.

\subsection{Basic and algebraic}

\begin{definition}[Basic, algebraic]\label{d:basic}
Let $I$ be a small discrete category, $R\in\Delta$ and $i\colon I\to R$ a functor. Then $i$ is \emph{basic} if there is a sequence $(R_\alpha)_{\alpha\le\gamma}$ in $\Delta$, indexed by a small ordinal $\gamma$, with $D$-functors $R_\alpha\to R_{\alpha'}$ for $\alpha\le\alpha'$, such that
\begin{itemize}
\item $R_0=\widetilde I$;
\item each $R_\alpha\to R_{\alpha+1}$ is a localisation, an equalisation or an adjunction of arrows, with small data;
\item for a limit ordinal $\lambda$, $R_\lambda$ is the 2-colimit in $\Delta$ of the $R_\beta$, $\beta<\lambda$;
\item there is an equivalence $R_\gamma\simeq R$ in $\Delta$ whose composite with $I\to\widetilde I\to R_\gamma$ is isomorphic to $i$.
\end{itemize}
The category $R$ is \emph{algebraic} if some functor $I'\to R$ from a small discrete category is basic.
\end{definition}

This is the meaning of \q{basic} on his board at \ts{4gynb9uTEE4}{2h0m40s}{2:00:40}, where the steps are called B (localisation and equalisation, that is, imposing axioms) and C (adjoining arrows, that is, adding data).

\subsection{Generation and saturation}

\begin{definition}[Condition (G), saturation]\label{d:G}
A functor $i\colon I\to R$ satisfies \emph{(G)} if every strictly full subcategory $R'\subseteq R$ that contains $i(I)$ and is closed in $R$ under $D'$-limits and $D''$-colimits is all of $R$. The \emph{saturation} of $i$ is the sequence of replete full subcategories $I_\alpha$ defined by: $I_0$ is the repletion of $i(I)$; $I_{\alpha+1}$ is $I_\alpha$ together with all the objects isomorphic to a limit of a $D'$-diagram in $I_\alpha$ or to a colimit of a $D''$-diagram in $I_\alpha$; and $I_\lambda=\bigcup_{\beta<\lambda}I_\beta$ for a limit ordinal $\lambda$.
\end{definition}

\begin{proposition}\label{p:sat}
Condition (G) holds if and only if $\bigcup_{\alpha}I_\alpha=R$, the union over all ordinals. If every shape in $D$ has fewer than $\rho$ arrows, for a regular cardinal $\rho$, then (G) holds if and only if $I_\rho=R$.
\end{proposition}

This is the construction and the bound of \cite[\S3.5]{kellybook}.

\begin{proof}
The union $\bigcup_\alpha I_\alpha$ is replete and contains $i(I)$. It is closed under $D$-limits and $D$-colimits, because a diagram has a small set of vertices, each in some $I_\alpha$, so the diagram lands in a single $I_\alpha$. By induction on $\alpha$, each $I_\alpha$ is contained in every $R'$ as in (G). So the union is the smallest such $R'$, and (G) says that it is $R$. For the bound: a diagram with fewer than $\rho$ arrows in $I_\rho$ has fewer than $\rho$ vertices, each in some $I_\alpha$ with $\alpha<\rho$; by regularity it lands in one of them, so its limit or colimit lies in $I_{\alpha+1}\subseteq I_\rho$. Hence $I_\rho$ is already closed.
\end{proof}

So for finitary $D$ at most $\omega$ steps are needed, and in general the number of steps can exceed any given small ordinal only when the shapes of $D$ are unbounded. One step need not suffice even for colimits: Kelly observes that $\Cat_0$ is the colimit closure of $\{2\}$, the arrow category, while the category with one non-identity idempotent is not a colimit of a diagram of copies of $2$ \cite[end of \S5.9]{kellybook}.

\begin{remark}[Faithfulness and density]\label{r:faith}
Condition (G) implies that restriction $\Hom_D(R,C)\to C^I$ is faithful for every $C\in\Delta$. Indeed, given two transformations $\sigma,\tau\colon\psi\to\psi'$ of $D$-functors that agree on $i(I)$, the objects $X$ with $\sigma_X=\tau_X$ form a replete full subcategory containing $i(I)$; it is closed under $D$-limits, because an arrow into the limit $\psi'(\lim d)=\lim\psi' d$ is determined by its components, and dually under $D$-colimits. So (G) lies between \q{basic} (Theorem~\ref{t:basicG}) and faithfulness, which he himself judged too weak. Condition (G) is also weaker than codensity in Isbell's sense \cite{isbell} (\q{right adequate}, \cite[\S5.1]{kellybook}), which asks every object to be the canonical limit of the objects of $I$ under it, and so gives (G) for limits in one step. In $\Set$, $\{2\}$ satisfies (G) for limits (Corollary~\ref{c:twosteps}) but is not codense: no set with three elements is even a limit of copies of $2$.
\end{remark}

\section{Basic implies (G), for every $D$}\label{sec:basicG}

\begin{lemma}\label{l:ff}
A functor $i\colon I\to R$ satisfies (G) if and only if, for every fully faithful $D$-functor $u\colon B'\to R$ with $B'\in\Delta$ and every functor $i'\colon I\to B'$ with $ui'\cong i$, the functor $u$ is an equivalence.
\end{lemma}

\begin{proof}
Assume the condition, and let $R'$ be as in (G). Since $R'$ is closed under $D$-limits and $D$-colimits computed in $R$, it lies in $\Delta$ and its inclusion $u$ is a $D$-functor; $i$ factors through $R'$ because $R'$ is replete and contains $i(I)$. So $u$ is an equivalence, and a replete full subcategory equivalent to $R$ by its inclusion is $R$. Conversely, the essential image of a functor $u$ as in the statement is replete and contains $i(I)$. It is closed under $D'$-limits: a $D'$-diagram on objects $u(b_k)$ lifts to $B'$ because $u$ is fully faithful, its limit exists in $B'$, and $u$ preserves it. Colimits are dual. So by (G) the essential image is $R$, and $u$ is an equivalence.
\end{proof}

\begin{theorem}\label{t:basicG}
Assume (H). If $i\colon I\to R$ is basic, then $i$ satisfies (G). This holds for every class $D$, small or not, and also for a small category $I$ that is not discrete.
\end{theorem}

\begin{proof}
Let $(R_\alpha)_{\alpha\le\gamma}$ be a basic sequence. We show by transfinite induction that $I\to R_\alpha$ satisfies (G) for every $\alpha$. At each stage, let $u\colon B'\to R_\alpha$ and $i'$ be as in Lemma~\ref{l:ff}; replacing $B'$ by its essential image, we may assume that $u$ is the inclusion of a replete full subcategory. We show that $u$ is essentially surjective; then it is an equivalence.

\emph{The envelope, $\alpha=0$.} We have $i'\colon I\to B'$ with $ui'\cong j$. By the universal property of $\widetilde I$, with $C=B'$, extend $i'$ to a $D$-functor $r\colon\widetilde I\to B'$ with $rj\cong i'$. Then $urj\cong ui'\cong j=\mathrm{id}\circ j$, and since restriction along $j$ is fully faithful, with $C=\widetilde I$, we get $ur\cong\mathrm{id}$. So $u$ is essentially surjective. (For $D''$ a class of colimits and $D'=\varnothing$, this is Kelly's description of the free cocompletion as a closure \cite[Thm~5.35]{kellybook}.)

\emph{A successor step.} Let $F\colon R_\alpha\to R_{\alpha+1}$ be the step. The full subcategory $F^{-1}(B')$ of $R_\alpha$ is replete, contains the image of $I$, and is closed under $D$-limits and $D$-colimits, because $F$ preserves them and $B'$ is closed under them. By induction it is all of $R_\alpha$, so $F=uF'$ for a functor $F'\colon R_\alpha\to B'$, which is a $D$-functor because $u$ is fully faithful and reflects the limits and colimits that it preserves.
\begin{itemize}
\item Localisation at $\Sigma$: the functor $F'$ inverts $\Sigma$, because $uF'=F$ does and the fully faithful $u$ reflects isomorphisms. So $F'\cong rF$ for some $D$-functor $r\colon R_{\alpha+1}\to B'$. Then $urF\cong F=\mathrm{id}\circ F$, and since restriction along $F$ is fully faithful, $ur\cong\mathrm{id}$.
\item Equalisation: the same argument, since the faithful $u$ reflects equalities of arrows.
\item Adjunction of arrows: let $\gamma_h\colon Fa_h\to Fb_h$ be the universal arrows. Since $u$ is full and faithful, $\gamma_h=u(\beta'_h)$ for unique $\beta'_h\colon F'a_h\to F'b_h$. The pair $(F',\beta')$ gives a $D$-functor $r\colon R_{\alpha+1}\to B'$ with $(rF,r\gamma)\cong(F',\beta')$. Then $(urF,ur\gamma)\cong(F,\gamma)$, which corresponds under the equivalence $\Hom_D(R_{\alpha+1},R_{\alpha+1})\simeq\{(\psi,\beta)\}$ to $ur\cong\mathrm{id}$.
\end{itemize}

\emph{A limit ordinal $\lambda$.} By induction each $F_\beta\colon R_\beta\to R_\lambda$ factors as $uF'_\beta$. The coherence isomorphisms between the $F_\beta$ lift uniquely through the fully faithful $u$, so the $F'_\beta$ form a pseudo-cocone and induce $r\colon R_\lambda\to B'$ with $rF_\beta\cong F'_\beta$ compatibly. Then $urF_\beta\cong F_\beta$ for all $\beta$, and $ur\cong\mathrm{id}$ by the universal property of the 2-colimit.

\emph{The end.} Condition (G) is invariant under an equivalence $R_\gamma\simeq R$ compatible with $I$.

In each case the universal property is used with target $B'$, which lies in $\Delta$; no smallness of $D$ is needed. Nothing used the discreteness of $I$.
\end{proof}

\section{Reduction of the converse to algebraicity}\label{sec:reduction}

\begin{lemma}\label{l:closure}
Assume (H). Let $i\colon I\to R$ be a functor from a small discrete category, and let $x$ be an object of the closure of $i(I)$ in $R$ under $D$-limits and $D$-colimits. Then there are a basic $I\to S$, a $D$-functor $\Phi\colon S\to R$ whose restriction to $I$ is isomorphic to $i$, and an object $s\in S$ with $\Phi s\cong x$. The same holds for a small family of objects $(x_k)$ of the closure, with one $S$ and objects $s_k$.
\end{lemma}

\begin{proof}
By induction on the saturation stage of $x$ (Proposition~\ref{p:sat}). If $x$ is in $I_0$ take $S=\widetilde I$ and $\Phi$ the extension of $i$. Otherwise let $x\cong\lim d$ with $d\colon J\to I_\beta$ a $D'$-diagram; the colimit case is dual.
\begin{enumerate}
\item For each vertex $d_k$, $k\in J$, choose by induction a triple $(S_k,\Phi_k,s_k)$ with $\Phi_ks_k\cong d_k$.
\item Form the amalgamated 2-sum of the $S_k$ over $\widetilde I$, an iterated 2-pushout, with the $D$-functor to $R$ induced by the $\Phi_k$. It is basic over $I$: the steps of each $S_k$, pushed out along $\widetilde I\to$ the sum, are steps of the same kind, as one sees on the descriptions of $\Hom_D(-,C)$ in Definition~\ref{d:H}. This construction appears in the course (\ts{PqCQM3Pcl6Q}{9h21m6s}{9:21:06}).
\item Adjoin arrows $s_k\to s_{k'}$, one for each arrow $k\to k'$ of $J$, mapped by $\Phi$ to the arrows of $d$ transported along the isomorphisms $\Phi_ks_k\cong d_k$. They form a small set, because $J$ is small, and $\Phi$ extends.
\item Equalise the composition and identity relations of $J$ among the new arrows. They hold after applying $\Phi$, so $\Phi$ extends again.
\item In the last stage $S$ there is now a diagram $\tilde d\colon J\to S$ with $\Phi\tilde d\cong d$. Put $s=\lim\tilde d$; then $\Phi s\cong\lim d\cong x$, since $\Phi$ preserves $D'$-limits.
\end{enumerate}
For a small family, apply step (2) to the triples of all the $x_k$ at once.
\end{proof}

\begin{theorem}\label{t:reduction}
Assume (H). A functor $i\colon I\to R$ from a small discrete category is basic if and only if it satisfies (G) and $R$ is algebraic.
\end{theorem}

\begin{proof}
If $i$ is basic, (G) holds by Theorem~\ref{t:basicG} and $R$ is algebraic by definition. Conversely, let $I'\to R$ be basic, through a sequence $(R_\alpha)_{\alpha\le\gamma}$ starting at $\widetilde{I'}$, and assume (G) for $i$. By Proposition~\ref{p:sat} every object $x\in I'$ lies in the closure of $i(I)$. Lemma~\ref{l:closure}, applied to the small family $(x)_{x\in I'}$, gives $(S,\Phi)$ basic over $I$, objects $s_x\in S$ and isomorphisms $c_x\colon x\cong\Phi s_x$ in $R$. We build from $R$, by small steps, a category $P^*$ basic over $I$ and equivalent to $R$.
\begin{enumerate}
\item Let $P=R\sqcup_{\widetilde{I'}}S$ be the 2-pushout of $R\leftarrow\widetilde{I'}\to S$, the second functor extending $x\mapsto s_x$, with legs $\iota_R$, $\iota_S$ and isomorphism $\theta\colon\iota_R|_{I'}\cong\iota_S(s_{-})$. It is basic over $I$: put $P_\alpha=R_\alpha\sqcup_{\widetilde{I'}}S$; then $P_0\simeq S$; the 2-pushout of a localisation, an equalisation or an adjunction of arrows is a step of the same kind, as one checks on the descriptions of $\Hom_D(-,C)$; and 2-pushouts commute with 2-colimits of chains.
\item We make the two functors $\iota_S$ and $\iota_R\Phi\colon S\rightrightarrows P$ isomorphic. Let $T_\beta\colon S_\beta\to S$, $\beta\le\delta$, be the functors of the basic sequence of $S$ over $I$. For $C\in\Delta$, let $M_\beta(C)$ be the category of pairs $(\chi,\varepsilon)$ with $\chi\in\Hom_D(P,C)$ and $\varepsilon$ an isomorphism $\chi\iota_ST_\beta\cong\chi\iota_R\Phi T_\beta$; a morphism is a transformation $\chi\to\chi'$ compatible with the $\varepsilon$'s. We build by small steps $D$-functors $Q_\beta\colon P\to P'_\beta$ and isomorphisms $\varepsilon_\beta$ that are universal: $\sigma\mapsto(\sigma Q_\beta,\sigma\varepsilon_\beta)$ is an equivalence $\Hom_D(P'_\beta,C)\simeq M_\beta(C)$ for every $C$, compatible with the restrictions $M_{\beta'}(C)\to M_\beta(C)$ for $\beta\le\beta'$.
\begin{itemize}
\item $\beta=0$, $S_0=\widetilde I$: adjoin arrows $e_k\colon\iota_S(k)\to\iota_R\Phi(k)$ and $e'_k$ in the other direction, for $k\in I$, and equalise $e'_ke_k$ and $e_ke'_k$ with the identities. This represents the isomorphisms between the restrictions to $I$, which are the isomorphisms between the functors on $\widetilde I$, since restriction $\Hom_D(\widetilde I,C)\to C^I$ is an equivalence.
\item A localisation or an equalisation $S_\beta\to S_{\beta+1}$: restriction $\Hom_D(S_{\beta+1},C)\to\Hom_D(S_\beta,C)$ is fully faithful, so $M_{\beta+1}(C)\to M_\beta(C)$ is an isomorphism, and $P'_{\beta+1}=P'_\beta$.
\item An adjunction of a set $H$ of arrows: by Definition~\ref{d:H}(4), an isomorphism between $D$-functors on $S_{\beta+1}$ is an isomorphism between their restrictions to $S_\beta$ that is compatible with the arrows $h\in H$. So $P'_{\beta+1}$ is $P'_\beta$ with the two sides of the naturality square of $\varepsilon_\beta$ at each $h$ equalised ($|H|$ pairs).
\item A limit ordinal $\lambda$: let $P'_\lambda$ be the 2-colimit of the $P'_\beta$, $\beta<\lambda$. As $S_\lambda$ is the 2-colimit of the $S_\beta$, an isomorphism between $D$-functors on $S_\lambda$ is a compatible family of isomorphisms between their restrictions to the $S_\beta$; so $M_\lambda(C)$ is equivalent to the pseudo-limit of the $M_\beta(C)$, $\beta<\lambda$. So is $\Hom_D(P'_\lambda,C)$, the pseudo-limit of the equivalent chain $\Hom_D(P'_\beta,C)$. Pseudo-limits of equivalent chains are equivalent, and $\varepsilon_\lambda$ is the image of $\mathrm{id}_{P'_\lambda}$.
\end{itemize}
Put $P'=P'_\delta$ and let $e$ be the isomorphism $Q_\delta\iota_S\cong Q_\delta\iota_R\Phi$ given by $\varepsilon_\delta$, which is defined on all of $S$ since $T_\delta$ is an equivalence.
\item For each $x\in I'$, equalise the two isomorphisms $\iota_R x\to\iota_R\Phi s_x$ given by $e_{s_x}\circ\theta_x$ and by $\iota_R(c_x)$ ($|I'|$ pairs). Call the result $P^*$.
\end{enumerate}
A $D$-functor $P^*\to C$ amounts to a tuple $(\psi,\varphi,\theta,e)$, with $\psi\colon R\to C$ and $\varphi\colon S\to C$ $D$-functors, $\theta$ an isomorphism between their restrictions to $I'$ and $e\colon\varphi\cong\psi\Phi$, subject to the equations of step (3). Such a tuple is determined by $\psi$ up to unique isomorphism: $\varphi\cong\psi\Phi$ through $e$, and then $\theta$ is forced by step (3). So restriction along the $R$-leg, $\Hom_D(P^*,C)\to\Hom_D(R,C)$, is an equivalence for every $C\in\Delta$. By the 2-Yoneda lemma, applied with $C=R$ and $C=P^*$, the $R$-leg $R\to P^*$ is an equivalence in $\Delta$; it is compatible with $I$ through the isomorphisms $e_k$. Every step used small data and the sequence has small length, so $i$ is basic.
\end{proof}

\begin{corollary}\label{c:reduction}
Let $\Delta$ satisfy (H). Problem three has a positive answer in $\Delta$ if and only if every $R\in\Delta$ that satisfies (G) for some functor from a small discrete category is algebraic.
\end{corollary}

\section{Small classes $D$}\label{sec:small}

\begin{theorem}\label{t:small}
Let $D$ be a small class, each shape of which has fewer than $\rho$ arrows, for a regular cardinal $\rho$. Assume (H), with the constructions applied to small categories and small data giving small categories, as Grothendieck assumes (\ts{4gynb9uTEE4}{2h8m14s}{2:08:14}). Then for every small discrete $I$ and every $i\colon I\to R$, the functor $i$ is basic if and only if it satisfies (G). The base is fixed: $I$ is not enlarged.
\end{theorem}

\begin{proof}[First proof, by Theorem~\ref{t:reduction}]
By Proposition~\ref{p:sat}, $R=I_\rho$. By induction each $I_\alpha$ has a small set of isomorphism classes of objects, since there is only a small set of diagrams of each of the small set of shapes in a category with a small set of objects \cite[\S3.5]{kellybook}; so $R$ is essentially small. A skeleton $R_s$ of $R$ is algebraic over $\mathrm{Ob}\,R_s$, as Grothendieck shows (\ts{PqCQM3Pcl6Q}{9h26m59s}{9:26:59}): starting from the envelope of $\mathrm{Ob}\,R_s$, adjoin the arrows of $R_s$, equalise the composition and identity relations, and invert the comparison maps of the small set of $D$-diagrams of $R_s$; by the 2-Yoneda lemma the result is equivalent to $R$. Theorem~\ref{t:reduction} concludes.
\end{proof}

\begin{proof}[Second proof, direct]
Build a chain $S_\beta$, $\beta\le\rho$, of small categories in $\Delta$ with $D$-functors $\Phi_\beta\colon S_\beta\to R$ extending $i$, starting from $S_0=\widetilde I$. Each successor step (a) adjoins a formal arrow $a\to b$ for every arrow $\Phi_\beta a\to\Phi_\beta b$ of $R$, for all objects $a,b$ of $S_\beta$, and (b) equalises every pair of parallel arrows that $\Phi_\beta$ identifies; both sets of data are small, because $S_\beta$ is small. Since $\rho$ is regular and the shapes are $\rho$-small, the 2-colimit $S_\rho$ can be computed in $\Cat$: diagrams and parallel pairs in $S_\rho$ come from earlier stages, and a functor out of $S_\rho$ is a $D$-functor exactly when all its restrictions are. So $\Phi_\rho$ is fully faithful: every arrow between images of objects of $S_\rho$ has been adjoined at the next stage, and every pair identified has been equalised. Its essential image is closed under $D$-limits and $D$-colimits, since diagrams lift along a fully faithful functor; by (G) it is everything. A fully faithful and essentially surjective $D$-functor is an equivalence in $\Delta$, so $i$ is basic.
\end{proof}

In particular finite limits with $I$ infinite but small, and $\kappa$-small limits, fall under Theorem~\ref{t:small}: there is no counterexample to problem three with small $D$.

\section{Grothendieck topoi}\label{sec:topos}

Here $\Delta$ is the 2-category of Grothendieck topoi with the functors that preserve finite limits and small colimits, that is, the inverse images of geometric morphisms, in the opposite direction. Its objects satisfy Giraud's exactness conditions and have a small generating set \cite[exp.~IV]{sga4}, and the free constructions of (H) are classifying topoi, 2-pushouts (\ts{ORG14LNb5gE}{1h23m42s}{1:23:42}) and subtopoi. Read in the direction of geometric morphisms, the 2-pushouts and 2-colimits of chains of (H) are pullbacks and limits of chains of topoi. That all these exist is standard topos theory, which we use without reproving it.

\subsection{Two statements of the course on topoi}
The course on topoi states that every object of $R_{\alpha+1}$ is covered by images of objects of $R_\alpha$ (\ts{ORG14LNb5gE}{1h51m23s}{1:51:23}), and, as a necessary condition, that if $I\to R$ is basic then the finite products of the objects of $I$ generate $R$ (\ts{ORG14LNb5gE}{1h59m7s}{1:59:07}).

\begin{proposition}\label{p:toposfalse}
Both statements are false as stated.
\end{proposition}

\begin{proof}
Take $R_\alpha=\Set$ and adjoin one arrow $1\to1\sqcup1$. The result is the classifying topos of decompositions of $1$ into two summands, $R_{\alpha+1}=\Set\times\Set$, and the image of $\Set$ is the diagonal $\{(n,n)\}$. The object $(1,0)$ receives a map from $(n,n)$ only when $n=0$, so it is not covered by images of objects of $R_\alpha$.

For the second statement take $I=\{x\}$. Start from the object classifier, the envelope of $I$; invert $x\to1$, which gives $\Set$; then adjoin $1\to1\sqcup1$. This makes $\{(1,1)\}\to\Set\times\Set$ basic. But the finite products of $(1,1)$ are just $(1,1)$, which does not generate: the two maps $(1,0)\rightrightarrows(1,0)\sqcup(1,0)$ are distinct and agree after composition with every map from $(1,1)$, there being none.
\end{proof}

Condition (G) does hold in this example, as Theorem~\ref{t:topos} below requires: pulling back $(\mathrm{inl},\mathrm{inr})\colon(1,1)\to(2,2)$ along $\mathrm{inl}\colon(1,1)\to(2,2)$ gives $(1,0)$. An example of the same kind lives in $\mathbf Z/2$-sets: with $\tau$ the involution and $Y=\mathbf Z/2\sqcup1$, there is no equivariant map $Y^n\to\mathbf Z/2$, so $\{Y\}$ does not generate under finite products, yet $\mathbf Z/2=Y\times_{Y/\tau}1$, and the theory of torsors makes $\{Y\}$ basic. The sufficient condition stated in the course (\ts{ORG14LNb5gE}{2h0m13s}{2:00:13}) is not affected.

\subsection{Closures in a topos}

\begin{lemma}\label{l:topos}
Let $R$ be a Grothendieck topos and $K$ a small set of objects of $R$. The smallest replete full subcategory $B\subseteq R$ containing $K$ and closed under finite limits and small colimits is a Grothendieck topos, and its inclusion preserves finite limits and small colimits.
\end{lemma}

\begin{proof}
Write $R=\mathrm{Sh}(C,J)$ with $C$ a small site, let $\mu\ge\aleph_0$ bound the number of arrows of $C$, and call a sheaf $X$ \emph{small} if $\sum_c|X(c)|<\kappa$, where $\kappa=(2^\theta)^+$ and $\theta\ge\mu$ is so large that the objects of $K$ are small. Then $\nu^\mu<\kappa$ for every cardinal $\nu<\kappa$, since $\nu^\mu\le(2^\theta)^{\theta}=2^\theta$. We use four facts.
\begin{enumerate}
\item[(a)] Finite limits of small sheaves are small, being computed pointwise.
\item[(b)] If $F$ is a presheaf with fewer than $\kappa$ elements in all, its sheafification $F^{++}$ is small: $F^+(c)$ is a colimit over the covering sieves on $c$, of which there are at most $2^\mu$, of sets $\mathrm{Match}(S,F)$ of cardinality at most $\nu^\mu$ with $\nu<\kappa$. Hence coproducts in $R$ of fewer than $\kappa$ small sheaves, and images of small sheaves, are small.
\item[(c)] Every element $x\in X(c)$ of a sheaf lies in a small subsheaf, the image of $a(yc)\to X$, by (b), where $a(yc)$ is the sheafified representable.
\item[(d)] If maps $s_l\to X$, $l\in L$, are jointly epimorphic and $W\subseteq X$ is a small subsheaf, then the images of the $s_l$ for fewer than $\kappa$ indices $l$ already cover $W$: for each element $w$ of $W$, choose a covering sieve on each member of which $w$ lifts to some $s_l$, and one such $l$ for each member.
\end{enumerate}
Let $S$ be the class of small objects of $B$, which is essentially a small set, and $B_S$ the class of objects $X\in R$ on which the maps from objects of $S$ are jointly epimorphic. We show that $B\cap B_S$ contains $K$ and is closed under small colimits and finite limits; then $B\subseteq B_S$ by minimality of $B$.

It contains $K\subseteq S$, and the terminal object, which lies in $B$ and is small. A colimit in $R$ of objects of $B\cap B_S$ lies in $B$, and the maps from $S$ to it are jointly epimorphic, since the legs of a colimit are. For pullbacks, let $X\to Z\leftarrow Y$ be maps in $B\cap B_S$, and $W$ a small subsheaf of $X\times_ZY$; by (c) such subsheaves cover $X\times_ZY$, so it suffices to cover each $W$ by maps from $S$. The images $W_X\subseteq X$ and $W_Y\subseteq Y$ of $W$ are small by (b). By (d) and (b) there are $s_X\to X$ and $s_Y\to Y$, each a coproduct of fewer than $\kappa$ objects of $S$, hence in $S$, whose images contain $W_X$ and $W_Y$. Then $s_X\times_Zs_Y$ lies in $B$ and is small by (a), so it lies in $S$, and its image in $X\times_ZY$ is $\operatorname{im}(s_X)\times_Z\operatorname{im}(s_Y)\supseteq W$, epimorphisms being stable under pullback in a topos. So $X\times_ZY\in B_S$.

Hence $S$ is a small generating set of $B$. Finite limits and small colimits in $B$ are computed in $R$, so in $B$ coproducts are disjoint and universal, equivalence relations are effective and quotients are universal. By Giraud's theorem \cite[exp.~IV]{sga4}, $B$ is a Grothendieck topos.
\end{proof}

\begin{theorem}\label{t:topos}
Let $\Delta$ be the 2-category of Grothendieck topoi. A functor $i\colon I\to R$ from a small discrete category to a topos is basic if and only if it satisfies (G).
\end{theorem}

\begin{proof}
Every Grothendieck topos is algebraic, as Grothendieck notes (\ts{ORG14LNb5gE}{1h6m58s}{1:06:58}). Indeed, for $R=\mathrm{Sh}(C,J)$ with $C$ small, start from the envelope of $\mathrm{Ob}\,C$, the classifying topos of $\mathrm{Ob}\,C$-indexed families of objects; adjoining the arrows of $C$ and equalising its composition and identity relations gives the classifying topos of $C$-diagrams. The flat $J$-continuous diagrams form a quotient of that geometric theory, whose classifying topos $R$ is a subtopos, and a subtopos is the localisation at a small set of monomorphisms. The proof of Theorem~\ref{t:reduction} uses only constructions available for topoi, namely classifying topoi, 2-pushouts and subtopoi, so (G) implies basic.

For the other direction, the proof of Theorem~\ref{t:basicG} applies its universal properties to a replete full subcategory $B'$ that must lie in $\Delta$, and an arbitrary $R'$ as in (G) need not be a topos. By Lemma~\ref{l:topos}, applied at each stage $R_\alpha$ to the image of $I$, one may take for $B'$ the closure of that image under finite limits and small colimits: it is a topos, its inclusion is a 1-cell of $\Delta$, and it is contained in every $R'$ as in (G). The induction of Theorem~\ref{t:basicG} then shows that this closure is all of $R_\alpha$ at every stage, so (G) holds.
\end{proof}

The repaired form of the statement of the course is therefore \q{basic if and only if (G)}, with (G) in place of the covering and generation conditions of Proposition~\ref{p:toposfalse}.

\section{All small limits: the converse fails}\label{sec:set}

In this section $\Delta$ is either the 2-category of complete categories with continuous functors ($D'$ all small shapes, $D''=\varnothing$) or, dually, that of cocomplete categories with cocontinuous functors. Duality exchanges them: $I\to R$ is basic, or satisfies (G), in one exactly when $I\to R^{\mathrm{op}}$ is so in the other.

\subsection{Sets are generated by $2$ in two steps}

\begin{proposition}\label{p:eq}
Every set $X$ is an equaliser of two maps between powers of $2=\{0,1\}$:
\[X\cong\mathrm{Eq}\bigl(f,g\colon2^{P(X)}\rightrightarrows2^{P(P(X))\sqcup P(X)}\bigr),\]
\[f(u)=\bigl(S\mapsto u(\textstyle\bigcap S),\ A\mapsto u(A^c)\bigr),\qquad g(u)=\bigl(S\mapsto\textstyle\bigwedge_{A\in S}u(A),\ A\mapsto\neg u(A)\bigr),\]
where $P$ is the power set, $A^c$ the complement and $\bigcap\varnothing=X$.
\end{proposition}

\begin{proof}
The equaliser consists of the maps $u\colon P(X)\to2$ that preserve arbitrary intersections and complements, that is, the complete Boolean homomorphisms. For $x\in X$ the evaluation $u_x(A)=[x\in A]$ is one, and distinct points give distinct $u_x$. Conversely, let $u$ be one, and $F=\{A:u(A)=1\}$. Then $X=\bigcap\varnothing\in F$, so $\varnothing=X^c\notin F$; and $F$ is closed under arbitrary intersections, so $\bigcap F\in F$ is nonempty. Pick $x\in\bigcap F$. If $A\in F$ then $x\in A$; if $x\in A$ then $x\notin A^c$, so $A^c\notin F$, hence $u(A^c)=0$ and $u(A)=1$. Hence $u=u_x$, and $x$ is unique since $\{x\}\in F$.
\end{proof}

This is a form of the classical duality between $\Set$ and complete atomic Boolean algebras, and it also follows from Paré's theorem that the contravariant power-set functor $\Set^{\mathrm{op}}\to\Set$ is monadic \cite{pare}: every algebra is a coequaliser of free ones. We give the explicit pair because the Lean check below uses it.

\begin{corollary}\label{c:twosteps}
For $D$ all small limits, $\{2\}\to\Set$ satisfies (G), and its saturation reaches $\Set$ in exactly two steps: $I_2=\Set$ and $I_1\neq\Set$.
\end{corollary}

\begin{proof}
Powers of $2$ are in $I_1$, and Proposition~\ref{p:eq} puts every set in $I_2$. No set with three elements is in $I_1$. A limit of a diagram of copies of $2$ is the set of compatible families, a subset of a product of copies of $2$. Every map $2\to2$ commutes with the ternary operation $x+y+z$ of $\mathbf Z/2$, so this subset is closed under that operation, taken coordinatewise. A three-element subset $\{a,b,c\}$ closed under it would contain $a+b+c$; if $a+b+c=a$ then $b+c=0$, so $b=c$, and the other two cases are the same.
\end{proof}

\begin{proposition}\label{p:pres}
In $\Set^{\mathrm{op}}$ the object $2$ is not $\lambda$-presentable for any regular cardinal $\lambda$.
\end{proposition}

\begin{proof}
The functor $\Set^{\mathrm{op}}(2,-)$ preserves $\lambda$-filtered colimits exactly when $P=\Set(-,2)$ sends $\lambda$-cofiltered limits of $\Set$ to colimits of sets. Take the chain of tails $A_\alpha=[\alpha,\lambda)$, $\alpha<\lambda$, with the inclusions; its limit is $\bigcap_\alpha A_\alpha=\varnothing$. Then $P(\varnothing)$ has one element, while $\varinjlim_\alpha P(A_\alpha)$ has at least two: the germs of $\varnothing$ and of the whole $A_\alpha$ differ, since every $A_\beta$ is nonempty.
\end{proof}

\subsection{Algebraic categories are locally presentable}

\begin{proposition}\label{p:sketch}
Let $\Delta$ be the cocomplete categories with cocontinuous functors, and let $S\in\Delta$ be algebraic. Then $S$ is equivalent to a full reflective subcategory $\Phi\text{-}\mathrm{Alg}$ of a presheaf category $[A^{\mathrm{op}},\Set]$ on a small category $A$, closed there under $\lambda$-filtered colimits for some regular cardinal $\lambda$.
\end{proposition}

\begin{proof}
We follow a basic sequence $(S_\beta)$ and show by induction that $S_\beta\simeq\Phi_\beta\text{-}\mathrm{Alg}$ for a small sketch $(A_\beta,\Phi_\beta)$, in Kelly's sense, with
\[\mathrm{Cocts}[\Phi_\beta\text{-}\mathrm{Alg},B]\simeq\Phi_\beta\text{-}\mathrm{Com}[A_\beta,B]\]
for every cocomplete $B$, pseudo-naturally in $B$; at each step the sketch receives a functor from the previous one, along which restriction of comodels corresponds to restriction along the step. Here $\Phi\text{-}\mathrm{Alg}$ is the reflective orthogonality class of a regulus in $[A^{\mathrm{op}},\Set]$ \cite[Thm~6.5, Thm~6.11]{kellybook}, and the displayed equivalence is \cite[Thm~5.56]{kellybook}: with $B$ cocomplete and $A$ small, left adjoints may be replaced by cocontinuous functors. The steps:
\begin{itemize}
\item the envelope of $I$: the sketch $(I,\varnothing)$, which gives $\Set^I$;
\item enlarging the sketch, to make a given small set of objects of $\Phi\text{-}\mathrm{Alg}$ available in it: with $K\colon A\to\Phi\text{-}\mathrm{Alg}$ the dense functor, pass to the closure $D_0$ of $K(A)$ under $\lambda$-small colimits, for $\lambda$ so large that $D_0$ contains these objects (each is a small colimit of objects of $K(A)$); it is small \cite[\S3.5]{kellybook}, and $\Phi\text{-}\mathrm{Alg}\simeq(D_0,\lambda\text{-colimits})\text{-}\mathrm{Alg}$ with the same comodels in every cocomplete $B$ \cite[Prop.~6.21, Cor.~6.22, Thm~6.23]{kellybook};
\item inverting an arrow $\sigma\colon X\to Y$: enlarge so that $\sigma$ lies in $D_0$, then add the cylinder saying that $Y$, with leg $\sigma$, is the colimit of the one-object diagram $X$;
\item equalising a pair $X\rightrightarrows Y$: enlarge so that the coequaliser $Y\to Q$ lies in $D_0$, and invert $Y\to Q$;
\item adjoining arrows: adjoin the arrows freely to $D_0$, keeping the same cylinders;
\item a limit ordinal $\lambda$: let $A_\lambda$ be the pseudo-colimit in $\Cat$ of the $A_\beta$, $\beta<\lambda$, which is small, and $\Phi_\lambda$ the images of the cylinders of all the $\Phi_\beta$. A functor $A_\lambda\to B$ is a compatible family of functors $A_\beta\to B$, and it is a comodel if and only if each of them is, so $\Phi_\lambda\text{-}\mathrm{Com}[A_\lambda,B]$ is the pseudo-limit of the $\Phi_\beta\text{-}\mathrm{Com}[A_\beta,B]$. By induction this chain is equivalent, pseudo-naturally in $B$, to the chain $\Hom_D(S_\beta,B)$, whose pseudo-limit is $\Hom_D(S_\lambda,B)$; and $\Phi_\lambda\text{-}\mathrm{Com}[A_\lambda,B]\simeq\mathrm{Cocts}[\Phi_\lambda\text{-}\mathrm{Alg},B]$ by the theorems just cited. By the 2-Yoneda lemma, $B$ ranging over $\Delta$, $S_\lambda\simeq\Phi_\lambda\text{-}\mathrm{Alg}$, compatibly with the $S_\beta$.
\end{itemize}
Finally, the cylinders of the last sketch are cones over diagrams of presheaves of size $<\lambda$ for some regular $\lambda$, so the orthogonality class is closed in $[A^{\mathrm{op}},\Set]$ under $\lambda$-filtered colimits.
\end{proof}

\begin{theorem}\label{t:notalg}
Let $\Delta$ be the cocomplete categories. Every object of an algebraic $S\in\Delta$ is $\mu$-presentable for some regular $\mu$. Hence $\Set^{\mathrm{op}}$ is not algebraic, and $\{2\}\to\Set^{\mathrm{op}}$ satisfies (G) without being basic. Dually, for $D$ all small limits, $\{2\}\to\Set$ satisfies (G) and is not basic, and $\Set$ is not algebraic.
\end{theorem}

This is the counterexample he pointed to on 10 July (\ts{4gynb9uTEE4}{3h43m0s}{3:43:00}, \ts{4gynb9uTEE4}{3h44m8s}{3:44:08}).

\begin{proof}
Take $S\simeq\Phi\text{-}\mathrm{Alg}$ as in Proposition~\ref{p:sketch}, closed under $\lambda$-filtered colimits in $[A^{\mathrm{op}},\Set]$. Every presheaf on the small category $A$ is a small colimit of representables, which are finitely presentable, so it is $\mu$-presentable for some regular $\mu\ge\lambda$. Since $\Phi\text{-}\mathrm{Alg}$ is a full subcategory closed under $\mu$-filtered colimits, an object of it that is $\mu$-presentable in $[A^{\mathrm{op}},\Set]$ is $\mu$-presentable in $\Phi\text{-}\mathrm{Alg}$. By Proposition~\ref{p:pres}, $\Set^{\mathrm{op}}$ is therefore not algebraic. By Corollary~\ref{c:twosteps} and duality, $\{2\}\to\Set^{\mathrm{op}}$ satisfies (G), so it is not basic, by definition of \q{algebraic}. The last sentence is the dual statement.
\end{proof}

Theorem~\ref{t:notalg} needs no part of (H): \q{basic} only asks that the steps used exist. That algebraic categories are, up to duality, the locally presentable ones is Gabriel--Ulmer duality \cite{gu}, \cite[ch.~1]{ar}, and $\Set^{\mathrm{op}}$ is not locally presentable \cite{gu}; the argument above uses only the results of \cite{kellybook} cited, and one half of that duality.

\begin{remark}[Lean]\label{r:lean}
Proposition~\ref{p:eq}, Corollary~\ref{c:twosteps} and Proposition~\ref{p:pres} are checked in Lean~4 with Mathlib, with the standard axioms only, in the file \texttt{lean/\allowbreak Buffalo/\allowbreak ProblemThree.lean} of the edition's repository \cite{repo} (commit \texttt{23657b5d}), with $\Set$ the types of a fixed universe. The declaration \texttt{equivEqualizer} is the bijection of every type with the equaliser of Proposition~\ref{p:eq}. The declarations \texttt{saturation\_two} and \texttt{not\_saturation\_one} say that the saturation of $\{2\}$ under limits of small diagrams has every type at its second stage and no three-element type at its first, through \texttt{Buffalo.Topos.Part05.no\_limit\_of\_bools}, the closure under the ternary operation; and \texttt{conditionG} is (G) for $\{2\}\to\Set$. The declarations \texttt{not\_isCardinalPresentable} and \texttt{not\_isPresentable} say that $2$ is not $\kappa$-presentable in $\Set^{\mathrm{op}}$ for any regular $\kappa$, with Mathlib's notion of presentability and the chain of tails of $\kappa$. Proposition~\ref{p:sketch} and Theorem~\ref{t:notalg} are not formalised: neither the notions \q{basic} and \q{algebraic} nor hypothesis (H) are modelled, nor Kelly's sketches, so the conclusions that $\Set^{\mathrm{op}}$ is not algebraic and that $\{2\}\to\Set$ is not basic rest on the argument above.
\end{remark}

\subsection{The number of steps}

\begin{proposition}\label{p:bound}
Let $\Delta$ be the cocomplete categories, let $R\in\Delta$ be algebraic, and let $i\colon I\to R$ satisfy (G), with $I$ small. Then $I_{\alpha+1}=R$ for some small ordinal $\alpha$. Dually for complete categories.
\end{proposition}

\begin{proof}
Write $R\simeq\Phi\text{-}\mathrm{Alg}$ as in Proposition~\ref{p:sketch}, with reflector $L$ and inclusion $J$. For every object $X$, the presheaf $JX$ is the colimit of representables $yA$ over its category of elements, and $L$ preserves colimits, so $X\cong LJX$ is a colimit in $R$ of objects $L(yA)$, $A\in A$. These form a small set $G$. By (G) and Proposition~\ref{p:sat}, each object of $G$ lies in some $I_{\alpha_g}$; let $\alpha$ be the supremum of the $\alpha_g$, a small ordinal since $G$ is small. Every object of $R$ is then a colimit of a small diagram in the full subcategory $I_\alpha$, so $I_{\alpha+1}=R$.
\end{proof}

Adámek and Rosický give a cocomplete, locally small category $\mathcal C$ of structures with partial unary operations indexed by all ordinals, and an object $A$ whose closure under colimits in $\mathcal C$ is reached in no small number of steps \cite[Example~6.38]{ar}. Rosický pointed to this example in answer to a question of Campion on MathOverflow, and Campion's own answer there works it out \cite{mo}. On the same page Barr recalls a claim of Isbell, which he believes was never published, that the limit closure of the group $\mathbf Z/2$, seen as a one-object category, also needs a proper class of steps.

\begin{corollary}\label{c:ar}
Let $R$ be the closure of $\{A\}$ under small colimits in the category $\mathcal C$ of \cite[Example~6.38]{ar}. Then $R$ is cocomplete, $\{A\}\to R$ satisfies (G), its saturation reaches $R$ in no small number of steps, and $R$ is not algebraic. So in the 2-category of cocomplete categories, $\{A\}\to R$ is a second counterexample to problem three, and no small ordinal bounds the number of steps for functors satisfying (G).
\end{corollary}

\begin{proof}
Colimits in $R$ are computed in $\mathcal C$, so $R$ is cocomplete and locally small, and (G) holds by construction; the saturation of $\{A\}$ in $R$ is its saturation in $\mathcal C$. By the example, $I_\alpha\neq R$ for every small $\alpha$. By Proposition~\ref{p:bound}, $R$ is not algebraic.
\end{proof}

Proposition~\ref{p:bound} and Corollary~\ref{c:twosteps} show that a bounded number of steps is necessary for algebraicity and not sufficient: $\{2\}\to\Set$ needs two steps and $\Set$ is not algebraic. Rosický also notes there that under Vop\v{e}nka's principle the colimit closure of a small full subcategory of a locally presentable category is locally presentable \cite{mo}.

\section{Open questions}\label{sec:open}

By Corollary~\ref{c:reduction}, for each $\Delta$ satisfying (H), problem three asks whether every $R\in\Delta$ that satisfies (G) for some small $I$ is algebraic. The answer is yes for small $D$ (Theorem~\ref{t:small}) and for topoi (Theorem~\ref{t:topos}), and no for all small limits and for all small colimits (Theorem~\ref{t:notalg}, Corollary~\ref{c:ar}).

\begin{question}\label{q:mixed}
Let $\Delta$ be the 2-category of categories with finite limits and all small colimits, with the functors preserving them, and no exactness condition. Does (G) for some small $I$ imply that $R$ is algebraic? In particular, $\{2\}\to\Set^{\mathrm{op}}$ satisfies (G) in this $\Delta$, since it does under colimits alone: is $\Set^{\mathrm{op}}$ algebraic in it?
\end{question}

For mixed large classes such as this one we do not even know whether (H) holds; the localisations and the 2-colimits of chains may fail to be locally small. Proposition~\ref{p:sketch} does not apply, because its proof uses sketches for colimits only.

\begin{question}\label{q:steps}
For mixed large classes $D$, is the number of saturation steps of a basic $I\to R$ bounded by a small ordinal?
\end{question}

For $D$ all small colimits, or all small limits, the answer is yes for basic functors (Proposition~\ref{p:bound}) and no for functors that only satisfy (G) (Corollary~\ref{c:ar}). Grothendieck expected a bound in general, without proof (\ts{PqCQM3Pcl6Q}{8h8m14s}{8:08:14}).

The second question of 2 July, whether the base of an algebraic $R$ can be enlarged so that only axioms remain to be imposed, has a positive answer in the locally presentable case by Gabriel--Ulmer: taking for $I$ the objects of the opposite of the category of $\lambda$-presentable objects, the category is that of $\lambda$-continuous functors on it, so that only axioms remain \cite{gu}. We have not examined it for the other classes.

\subsection*{What is not done}
Hypothesis (H) is assumed, as Grothendieck assumed it, and not proved, except where Kelly's book provides it (Section~\ref{sec:framework}); in particular we have no reference for mixed small classes. For topoi, the existence of the constructions (classifying topoi, subtopoi, pullbacks and limits of chains of topoi) is taken from standard topos theory without proof. Of the counterexample for all small limits, the statements about sets are formalised in Lean (Remark~\ref{r:lean}); the notions \q{basic} and \q{algebraic}, the free constructions and Kelly's sketches are not, so Proposition~\ref{p:sketch}, Theorem~\ref{t:notalg} and the positive results rest on the proofs given here. Example~6.38 of Adámek and Rosický is used as described on MathOverflow by Rosický and Campion \cite{mo}; we have not had access to the book itself. None of the arguments has been checked by a specialist.

\subsection*{Use of language models}
The question was located in the recordings of the Buffalo course, and the literature searched, with language models (Claude Opus~5.5, Anthropic), which also found the arguments of this note, wrote the Lean file and a first draft of this text, under the author's direction. The Lean statements are checked by the Lean kernel; the other proofs have not been checked by a specialist. The references were opened where an open copy exists, Kelly's book among them; what was checked for each one is listed in \texttt{docs/problem-three/review/citations.md} of the repository \cite{repo}.

\begin{thebibliography}{9}
\bibitem{ar} J.~Adámek, J.~Rosický, \emph{Locally Presentable and Accessible Categories}, LMS Lecture Note Ser.~189, Cambridge University Press, 1994, \href{https://doi.org/10.1017/CBO9780511600579}{doi:10.1017/CBO9780511600579}.
\bibitem{gu} P.~Gabriel, F.~Ulmer, \emph{Lokal präsentierbare Kategorien}, LNM~221, Springer, 1971, \href{https://doi.org/10.1007/BFb0059396}{doi:10.1007/BFb0059396}.
\bibitem{sga4} M.~Artin, A.~Grothendieck, J.-L.~Verdier (eds.), \emph{Théorie des topos et cohomologie étale des schémas} (SGA~4), tome~1, LNM~269, Springer, 1972, \href{https://doi.org/10.1007/BFb0081551}{doi:10.1007/BFb0081551}.
\bibitem{isbell} J.~R.~Isbell, Adequate subcategories, \emph{Illinois J.\ Math.} 4 (1960), 541--552, \href{https://doi.org/10.1215/ijm/1255456274}{doi:10.1215/ijm/1255456274}.
\bibitem{kellybook} G.~M.~Kelly, \emph{Basic Concepts of Enriched Category Theory}, LMS Lecture Note Ser.~64, Cambridge University Press, 1982; reprinted in \emph{Repr.\ Theory Appl.\ Categ.} 10 (2005), \url{http://www.tac.mta.ca/tac/reprints/articles/10/tr10abs.html}.
\bibitem{mo} T.~Campion, What's an example of a subcategory whose closure under colimits takes a lot of steps to form by iteration?, MathOverflow question 243931 (2016), with answers by J.~Rosický, T.~Campion and M.~Barr, \url{https://mathoverflow.net/questions/243931}.
\bibitem{pare} R.~Paré, Colimits in topoi, \emph{Bull.\ Amer.\ Math.\ Soc.} 80 (1974), 556--562, \href{https://doi.org/10.1090/S0002-9904-1974-13497-X}{doi:10.1090/S0002-9904-1974-13497-X}.
\bibitem{repo} Commutator-IO, grothendieck-project, \url{https://github.com/Commutator-IO/grothendieck-project}: index of results of the Buffalo topos course (\texttt{docs/buffalo-topos/index.tex}), the file \texttt{lean/Buffalo/ProblemThree.lean}, and the Buffalo tab of \url{https://grothendieck.commutator.io}.
\bibitem{open} M.~Hua, The open problems of Grothendieck's 1973 Buffalo lectures, and what became of them, 2026, \url{https://grothendieck.commutator.io/article/grothendieck-buffalo-open.pdf}.
\end{thebibliography}
\end{document}
